% not printed: handwritten call number QA153 W35 v.2; library stamps MATH.-STAT. LIBRARY and M777594
\begin{center}
{\large VORWORT ZUM ZWEITEN BANDE.}

\vspace{8pt}
\rule{2.5cm}{0.4pt}
\end{center}

\vspace{6pt}

Der in dem Vorworte zum ersten Bande angekündigten Absicht gemäss kann ich heute den zweiten Band meines Lehrbuches der Algebra der Oeffentlichkeit übergeben. Der dort aufgestellte Plan ist in den wesentlichen Punkten durchgeführt. Bei den Anwendungen bin ich bemüht gewesen, solche Probleme auszuwählen, die bereits in anderen Gebieten, der Geometrie oder Functionentheorie, ein selbständiges Interesse gewonnen haben, und die zugleich die Hauptpunkte der algebraischen Theorie möglichst vielseitig zur Anschauung bringen.

Die Anwendung der Theorie der algebraischen Zahlen ist bis zur Theorie der Kreistheilungszahlen durchgeführt. Wenn Leben und Arbeitskraft vorhalten, hoffe ich, in einer Fortsetzung meines Werkes die weiteren Anwendungen auf das Gebiet der elliptischen Functionen darzustellen, die nur zum Theil in meinem Buche „Elliptische Functionen und algebraische Zahlen“ enthalten sind.

Auch während der Ausarbeitung und des Druckes des zweiten Bandes hat mir die Hülfe und der Rath der Freunde zur Seite gestanden, die ich schon in der Vorrede zur ersten Auflage genannt habe. Aber auch manchen neuen Freund hat sich der erste Band bereits erworben, der meine Arbeit durch Winke und Rathschläge gefördert hat. Ihnen allen spreche ich an dieser Stelle meinen Dank aus, und füge die Bitte hinzu, dass sie dem Werke auch weiterhin ihr Interesse bewahren mögen.

\bigskip

\gesperrt{Strassburg}, im Juli 1896.

\begin{flushright}
\textbf{\gesperrt{Der Verfasser}.}\hspace*{2em}
\end{flushright}
